\documentclass[12pt]{article}
\usepackage[T1]{fontenc}
\usepackage[utf8]{inputenc}
\usepackage{lmodern}
\usepackage{textcomp}
\usepackage{enumitem}
\usepackage{geometry}
\geometry{margin=1in}
\begin{document}
\begin{center}
Answer Key - Astronomy - Graduate Level Assessment 18 \
\small (Answers are ordered exactly as in the question paper.)
\end{center}

\section*{Multiple Choice}
\begin{enumerate}[label=\arabic*.]
\item \textbf{Answer:} D
\\ \textit{Options:} A) 2:3 periodic ratio of Neptune:Pluto; B) Kirkwood Gaps.; C) Gaps in Saturn's rings.; D) Breaking of small Jovian moons to form ring materials.
\\ \textit{Note:} The breaking of small Jovian moons to form ring materials is primarily a result of tidal forces and collisions rather than orbital resonance, which typically influences the stability and motion of celestial bodies in relation to one another.
\item \textbf{Answer:} D
\\ \textit{Options:} A) The corona is made up mostly of charged particles which aren't luminous.; B) It's much too cool to emit visible light; C) We can't see magnetic fields; D) It's too diffuse
\\ \textit{Note:} The Sun's corona is too diffuse and faint compared to the intense brightness of the Sun's surface, making it visible only when the Sun is obscured during a total solar eclipse.
\item \textbf{Answer:} B
\\ \textit{Options:} A) pieces of comets rather than of asteroids.; B) chunks of large differentiated asteroids that were shattered by collisions.; C) chunks of rock chipped off the planet Mars.; D) leftover chunks of rock from the earliest period in the formation of the solar system.
\\ \textit{Note:} Meteorites with high metal content are likely fragments of differentiated asteroids, which have undergone processes that segregate their metallic core from silicate materials, and these asteroids were subsequently shattered by collisions in space.
\item \textbf{Answer:} C
\\ \textit{Options:} A) a New Horizons flyby in the 1990 s; B) Hubble Space Telescope images that resolved Pluto's disk; C) brightness measurements made during mutual eclipses of Pluto and Charon; D) radar observations made by the Arecibo telescope
\\ \textit{Note:} The correct answer is based on the fact that brightness measurements during mutual eclipses of Pluto and Charon allow astronomers to determine the size of Pluto by analyzing the changes in light as one body passes in front of the other, providing accurate data on their relative sizes and distances.
\item \textbf{Answer:} C
\\ \textit{Options:} A) It has a fusion crust; B) It contains solidified spherical droplets; C) It is highly processed; D) It has different elemental composition than earth
\\ \textit{Note:} A highly processed rock suggests extensive alteration or refinement by human activity or geological processes on Earth, which is inconsistent with the natural formation of meteorites that typically retain their original extraterrestrial characteristics.
\end{enumerate}

\section*{True / False}
\begin{enumerate}[label=\arabic*.]
\setcounter{enumi}{5}
\item \textbf{Answer:} False
\\ \textit{Options:} A) True; B) False
\\ \textit{Note:} The statement is false because weight is determined by mass and the acceleration due to gravity, which means that regardless of the materials' densities, five kilograms of feathers will always weigh more than one kilogram of bricks on the Moon or any other celestial body.
\item \textbf{Answer:} True
\\ \textit{Options:} A) True; B) False
\\ \textit{Note:} The relationship between actual brightness (luminosity) and apparent brightness (how bright an object appears from Earth) is governed by the inverse square law of light, which allows astronomers to calculate the distance to celestial objects when both values are known.
\item \textbf{Answer:} True
\\ \textit{Options:} A) True; B) False
\\ \textit{Note:} The tilt of the Earth's axis at approximately 23.5 degrees relative to its orbital plane, known as axial tilt or obliquity, is a fundamental characteristic of the Earth's orientation in space, significantly influencing seasonal variations and climate patterns.
\item \textbf{Answer:} True
\\ \textit{Options:} A) True; B) False
\\ \textit{Note:} Cassiopeia is classified as a prominent constellation in the northern sky due to its distinctive W shape formed by five bright stars, making it easily identifiable.
\item \textbf{Answer:} True
\\ \textit{Options:} A) True; B) False
\\ \textit{Note:} Dark energy is a hypothesized form of energy that permeates space and drives the observed acceleration of the universe's expansion, as evidenced by measurements of distant supernovae and the cosmic microwave background radiation.
\end{enumerate}

\section*{Long Form Response}
\begin{enumerate}[label=\arabic*.]
\setcounter{enumi}{10}
\item \textbf{Answer:} When the temperature of a blackbody is reduced to half of its original value, the effects on its spectral characteristics are significant. According to the Stefan-Boltzmann law, the total power emitted by a blackbody is proportional to the fourth power of its temperature, resulting in a reduction of emitted power to 1/16 of its previous value. Additionally, Wien's displacement law indicates that the peak emission wavelength of a blackbody is inversely proportional to its temperature. Consequently, halving the temperature causes the peak emission wavelength to double, resulting in a shift towards longer wavelengths in the spectrum. These changes reflect the fundamental principles governing blackbody radiation and highlight how thermal energy influences emission properties.
\\ \textit{Note:} correct because it accurately applies the Stefan-Boltzmann law to demonstrate that reducing the temperature of a blackbody to half results in a power emission decrease to 1/16, and utilizes Wien's displacement law to explain that the peak emission wavelength shifts to a longer wavelength, illustrating the fundamental relationship between temperature, emitted power, and spectral characteristics of blackbodies.
\item \textbf{Answer:} Constellations serve as a navigational and cultural framework for understanding the night sky, with their significance deeply rooted in human history and astronomy. The Milky Way, our galactic home, is a vast, disk-shaped structure composed of stars, gas, and dust, and its plane is defined by the distribution of these components. The constellation Leo, which represents a lion in mythology, is situated in such a way that it does not lie along the Milky Way's plane; instead, it is positioned above it. This occurs because Leo is located in the galactic halo, an area populated by older stars that are not confined to the plane of the galaxy. Understanding the positioning of constellations like Leo provides insight into the complex structure of the Milky Way and the distribution of stellar populations within and outside its disk.
\\ \textit{Note:} correct because it outlines the cultural and navigational importance of constellations like Leo while clarifying its physical position relative to the Milky Way's plane, which is determined by the galactic structure of stars, gas, and dust.
\end{enumerate}

\vfill
\noindent\textit{End of Answer Key}
\end{document}